\documentclass[letterpaper,10pt]{article}

\usepackage{latexsym}
\usepackage[empty]{fullpage}
\usepackage{titlesec}
\usepackage{marvosym}
\usepackage[usenames,dvipsnames]{color}
\usepackage{verbatim}
\usepackage{enumitem}
\usepackage[hidelinks]{hyperref}
\usepackage{fancyhdr}
\usepackage[english]{babel}
\usepackage{tabularx}
\usepackage{fontawesome5}
\input{glyphtounicode}

\usepackage[sfdefault]{roboto}

\pagestyle{fancy}
\fancyhf{}
\fancyfoot{}
\renewcommand{\headrulewidth}{0pt}
\renewcommand{\footrulewidth}{0pt}

\addtolength{\oddsidemargin}{-0.5in}
\addtolength{\evensidemargin}{-0.5in}
\addtolength{\textwidth}{1in}
\addtolength{\topmargin}{-.5in}
\addtolength{\textheight}{1.0in}

\urlstyle{same}
\raggedbottom
\raggedright
\setlength{\tabcolsep}{0in}

\titleformat{\section}{\Large\bfseries\scshape\raggedright}{}{0em}{}[\titlerule]
\pdfgentounicode=1

\definecolor{tagbg}{RGB}{236,236,236}
\newcommand{\contacttag}[2]{%
  \colorbox{tagbg}{\strut\hspace{6pt}#1\hspace{4pt}#2\hspace{6pt}}%
}

\newcommand{\resumeItem}[1]{\item\small{#1}}
\newcommand{\resumeSubheading}[4]{
\vspace{-1pt}\item
  \begin{tabular*}{0.97\textwidth}[t]{l@{\extracolsep{\fill}}r}
    \textbf{#1} & #2 \\
    \textit{#3} & \textit{#4} \\
  \end{tabular*}\vspace{-7pt}
}
\newcommand{\resumeProjectHeading}[2]{
\item
  \begin{tabular*}{0.97\textwidth}{l@{\extracolsep{\fill}}r}
    \small\textbf{#1} & \small #2 \\
  \end{tabular*}\vspace{-7pt}
}
\newcommand{\resumeSubHeadingList}{\begin{itemize}[leftmargin=0.15in, label={}]}
\newcommand{\resumeSubHeadingListEnd}{\end{itemize}}
\newcommand{\resumeItemListStart}{\begin{itemize}[leftmargin=0.25in]}
\newcommand{\resumeItemListEnd}{\end{itemize}\vspace{-6pt}}

\begin{document}

\begin{center}
  \textbf{\Huge Devanshu Kumar Ranjan} \\
  \vspace{4pt}
  \footnotesize
  \contacttag{\faIcon{phone}}{+91-6206938457}\hspace{3pt}
  \contacttag{\faIcon{envelope}}{\href{mailto:devanshukranjan@gmail.com}{devanshukranjan@gmail.com}}\hspace{3pt}
  \contacttag{\faIcon{map-marker-alt}}{Patna, India}\hspace{3pt}
  \contacttag{\faIcon{linkedin}}{\href{https://linkedin.com/in/devanshu-kumar-ranjan-06399221b}{linkedin}}\hspace{3pt}
  \contacttag{\faIcon{github}}{\href{https://github.com/devanshukranjan}{github}}\hspace{3pt}
  \contacttag{\faIcon{globe}}{\href{https://devanshukranjan.github.io/Cybersecurity-Portfolio/}{portfolio}}
\end{center}

\section{Summary}
\small
Cybersecurity Engineer and AI Specialist with an M.Tech in Cyber Security from NIT Patna (CGPA: 9.06/10). Experienced in SOC operations, SIEM detection engineering, network traffic forensics, and applied machine learning for threat detection. Built a multi-VM Splunk SIEM enterprise lab simulating brute-force and web attack vectors with custom SPL hunting queries and real-time alerting. Analyzed packet-level DoS attack signatures with Wireshark and conducted digital forensic investigations using Autopsy and FTK Imager. Author of master's thesis on locally differentially private (LDP) sequential recommendation, securing user interaction data against inference attacks. Currently building production diagnostic engines and secure, high-throughput AI data systems at KAS Global Commerce. Holder of the Google Cybersecurity Specialization.
\normalsize

\section{Education}
\resumeSubHeadingList
  \resumeSubheading
      {National Institute of Technology, Patna}{Patna, India}
      {M.Tech., Cyber Security (CSE) | CGPA: 9.06/10}{August 2024 -- May 2026}
  \resumeSubheading
      {Graphic Era University, Dehradun}{Dehradun, India}
      {B.Tech., Computer Science and Engineering | CGPA: 8.83/10}{June 2019 -- June 2023}
  \resumeSubheading
      {Sri Chaitanya Techno School (CBSE)}{Visakhapatnam, India}{Senior Secondary | 80.6\%}{2019}
  \resumeSubheading
      {Cambridge School (ICSE)}{Cuttack, Odisha}{Secondary | 90.17\%}{2017}
\resumeSubHeadingListEnd

\section{Experience}
\resumeSubHeadingList
  \resumeSubheading
      {\href{https://mysellercentral.com/}{KAS Global Commerce Inc. (My Seller Central)}}{On-site}
      {AI Engineer}{July 2026 -- Present}
      \resumeItemListStart
          \resumeItem{\textbf{Root Cause Analysis (RCA) Engine (sole owner, research to production):} Architected an enterprise AI diagnostic system that analyzes anomalies across advertising performance datasets. Engineered statistical outlier algorithms ($\mu + 1.5\sigma$) and minimum-data threshold guards across 7 live diagnostics and 50+ root-cause rules.}
          \resumeItem{Implemented automated negative keyword recommendations with over-blocking protection to prevent wasted ad spend and mitigate budget exhaustion attacks.}
          \resumeItem{\textbf{Secure Data Architecture:} Guided teammates in maintaining resilient distributed scraping and multimodal insight generation pipelines across Amazon, Walmart, eBay, and Flipkart, ensuring access controls and rate-limiting resilience.}
      \resumeItemListEnd

  \resumeSubheading
      {\href{https://mysellercentral.com/}{KAS Global Commerce Inc. (My Seller Central)}}{On-site}
      {AI Research Intern}{September 2025 -- June 2026}
      \resumeItemListStart
          \resumeItem{Engineered high-throughput scraping infrastructure across Amazon, Walmart, and ONDC; reduced collection turnaround from 2--3 months to \textbf{1.5 weeks} via threaded Zyte API calls at \textbf{1,000 req/min} with secure MongoDB storage.}
          \resumeItem{Built taxonomy embedding pipelines using OpenAI \texttt{text-embedding-3-large} and PostgreSQL with \texttt{pgvector} for semantic category mapping, enforcing strict input sanitization and schema isolation.}
          \resumeItem{Designed category-level insight generation combining quantitative analytics with Gemini Flash inference across 25 best-selling competitor products per category.}
      \resumeItemListEnd
\resumeSubHeadingListEnd

\section{Projects}
\resumeSubHeadingList
    \resumeProjectHeading
        {Splunk SIEM Lab: Real Attack Detection \& Security Monitoring}{April 2026 -- May 2026}
        \resumeItemListStart
            \resumeItem{\textbf{Tools:} Splunk Enterprise, Splunk Universal Forwarder, Kali Linux, Ubuntu Server, Apache2, Hydra, Nikto, VMware Workstation}
            \resumeItem{Engineered a 3-VM virtualized SOC environment: Ubuntu Server with Splunk Universal Forwarder, Kali Linux attacker machine, and a dedicated Splunk Enterprise indexer/search head.}
            \resumeItem{Simulated SSH brute-force attacks with Hydra (600+ failed attempts) and web vulnerability scanning with Nikto, streaming attack telemetry into Splunk in real time.}
            \resumeItem{Formulated custom SPL hunting queries identifying brute-force bursts, top attacking source IPs, and 404 response anomalies; configured automated threshold-based alerts.}
            \resumeItem{Constructed a multi-panel Security Monitoring Dashboard visualizing attack timelines, targeted endpoints, and IP geographical distributions.}
            \resumeItem{GitHub: \href{https://github.com/devanshukranjan/Cybersecurity-Portfolio/tree/main/siem-investigations/splunk-siem-lab}{github.com/devanshukranjan/Cybersecurity-Portfolio/tree/main/siem-investigations/splunk-siem-lab}}
        \resumeItemListEnd

    \resumeProjectHeading
        {Privacy-Preserving Federated Sequential Recommendation (M.Tech Dissertation)}{August 2025 -- May 2026}
        \resumeItemListStart
            \resumeItem{\textbf{Tools:} Python, PyTorch, Differential Privacy, Federated Learning, Contrastive Learning}
            \resumeItem{Formulated \textbf{LDP-SeqMask} (Locally Differentially Private Sequence Masking), injecting calibrated noise into interaction sequences locally so raw client item sequences are never shared.}
            \resumeItem{Replaced unstable Exponential Mechanism item selection with a \textbf{Laplace perturbation mechanism}, eliminating gradient explosion and stabilizing federated updates.}
            \resumeItem{Applied \textbf{View-based Contrastive Learning (View-CL)} to maintain high recommendation accuracy under differential privacy noise, and integrated sequence-aware heterogeneous client model selection.}
        \resumeItemListEnd

    \resumeProjectHeading
        {Autonomous AI-Based Threat Detection and Elimination Engine}{May 2025 -- June 2025}
        \resumeItemListStart
            \resumeItem{\textbf{Tools:} Python, Tkinter, CNN, PyTorch, Scikit-learn, Pandas}
            \resumeItem{Engineered an AI threat detection system to detect ransomware activity and zero-day intrusion signatures.}
            \resumeItem{Trained and evaluated deep neural networks (CNNs) and ensemble models on the \textbf{UGRansome 2024} and \textbf{UNSW-NB15} benchmark datasets.}
            \resumeItem{Developed an interactive GUI-driven monitoring platform simulating live network traffic ingestion and automated threat containment.}
            \resumeItem{GitHub: \href{https://github.com/devanshukranjan/Autonomous-AI-Based-Threat-Detection-and-Threat-Elimination-Engine.git}{github.com/devanshukranjan/Autonomous-AI-Based-Threat-Detection-and-Threat-Elimination-Engine.git}}
        \resumeItemListEnd

    \resumeProjectHeading
        {DoS Attack Simulation \& Traffic Analysis}{January 2025 -- January 2025}
        \resumeItemListStart
            \resumeItem{\textbf{Tools:} Kali Linux, Metasploitable 2, Wireshark, hping3, Nmap, VirtualBox}
            \resumeItem{Executed 6 distinct DoS attack vectors (Slowloris, TCP SYN Flood, UDP Flood, TCP FIN Flood, PUSH+ACK Flood, ICMP Flood) against a Metasploitable target in a bridged virtual network.}
            \resumeItem{Captured and analyzed raw PCAP traffic in Wireshark, documenting packet inter-arrival patterns, TCP handshakes, and protocol-level anomaly signatures.}
            \resumeItem{GitHub: \href{https://github.com/devanshukranjan/Cybersecurity-Portfolio/tree/main/network-analysis/dos-attack-simulation}{github.com/devanshukranjan/Cybersecurity-Portfolio/tree/main/network-analysis/dos-attack-simulation}}
        \resumeItemListEnd

    \resumeProjectHeading
        {Digital Forensic Investigation \& File Recovery (Autopsy + FTK Imager)}{January 2025 -- February 2025}
        \resumeItemListStart
            \resumeItem{\textbf{Tools:} Autopsy, FTK Imager, Kali Linux, Raw Disk Image (.dd)}
            \resumeItem{Acquired forensic disk images in \texttt{.dd} format using FTK Imager, verifying cryptographic hashes (MD5 and SHA-1) to ensure evidence integrity.}
            \resumeItem{Conducted incident timeline reconstruction and artifact analysis using Autopsy, inspecting file system metadata, deleted structures, and system logs.}
            \resumeItem{Carved deleted files from unallocated clusters and verified file integrity using hex header and footer signature analysis (\texttt{\%PDF-} to \texttt{\%\%EOF}).}
            \resumeItem{GitHub: \href{https://github.com/devanshukranjan/Cybersecurity-Portfolio/tree/main/digital-forensics}{github.com/devanshukranjan/Cybersecurity-Portfolio/tree/main/digital-forensics}}
        \resumeItemListEnd
\resumeSubHeadingListEnd

\section{Technical Skills}
\resumeSubHeadingList
  \resumeItem{\textbf{SIEM \& SOC Operations}: Splunk Enterprise, Splunk Universal Forwarder, SPL Querying, Alerting, SOC Dashboards, Incident Response, Log Analysis}
  \resumeItem{\textbf{Threat Detection \& Network Security}: Wireshark, Nmap, Kali Linux, Metasploit, hping3, Hydra, Nikto, Slowloris, TCP/IP, IDS/IPS, Packet Analysis}
  \resumeItem{\textbf{Digital Forensics}: Autopsy, FTK Imager, File Carving, Timeline Reconstruction, Cryptographic Hash Verification (MD5/SHA-1), Hex Analysis}
  \resumeItem{\textbf{AI \& Security Intersections}: Locally Differential Privacy (LDP), Federated Learning, Adversarial Anomaly Detection, PyTorch, Scikit-learn, XGBoost}
  \resumeItem{\textbf{Languages \& Systems}: Python, Bash, C++, SQL, PostgreSQL, MongoDB, Linux / Ubuntu Server, VMware Workstation, VirtualBox, Git}
\resumeSubHeadingListEnd

\section{Certifications}
\resumeSubHeadingList
  \resumeItem{\href{https://www.coursera.org/account/accomplishments/specialization/QQOKJ9WNIWB9}{Google Cybersecurity Specialization (Coursera)} \hfill December 2024 -- March 2025}
  \resumeItem{\href{https://drive.google.com/file/d/1ofpmj-bLuC3EHXQGxpZV1dIxlXbkgEd6/view?usp=sharing}{Certified Python Developer Associate (PDAC)} \hfill March 2024}
  \resumeItem{\href{https://drive.google.com/file/d/1G3WGzaN9tsonASGrbLfKoXXj_7vUHSiA/view?usp=sharing}{Data Science for Engineers (NPTEL)} \hfill January 2022 -- March 2022}
  \resumeItem{\href{https://drive.google.com/file/d/19hY_be_Ezg3w1nLbYdGw9-6Z1nsBYbDd/view?usp=sharing}{Machine Learning (Teachnook)} \hfill February 2022 -- March 2022}
\resumeSubHeadingListEnd

\end{document}